\begin{blocksection}
\question Implement \texttt{compose}, which takes in two \lstinline{UserFn}s \texttt{f} and \texttt{g}, and returns a function \texttt{h} such that applying \texttt{h} to some arguments gives the same result as applying \texttt{g} to those arguments and then applying \texttt{f} to that result. For example, composing \lstinline|fn(x) { (x * x) }| with \lstinline|fn(x, y) { (y - x) }| gives \lstinline|fn(x, y) { ((y - x) * (y - x)) }|. Since \texttt{f} receives exactly one value, return a \lstinline{WrongArgumentCount} error if \texttt{f} does not take exactly one parameter. You may use \texttt{substitute}.
\begin{lstlisting}
pub fn compose(f: UserFn, g: UserFn) -> Result(UserFn, EvalError) {
  case f.params {
    ____________ -> Ok(________________________________________)
    params -> Error(WrongArgumentCount(want: ____________, got: 1))
  }
}

assert compose(
    UserFn(["x"], Binary(Mul, Var("x"), Var("x"))),
    UserFn(["x", "y"], Binary(Sub, Var("y"), Var("x"))),
  )
  == Ok(UserFn(
    ["x", "y"],
    Binary(Mul, Binary(Sub, Var("y"), Var("x")), Binary(Sub, Var("y"), Var("x"))),
  ))
assert compose(
    UserFn(["x"], Binary(Mul, Var("x"), N(2))),
    UserFn(["x"], Binary(Add, Var("x"), N(3))),
  )
  == Ok(UserFn(["x"], Binary(Mul, Binary(Add, Var("x"), N(3)), N(2))))
assert compose(UserFn(["x", "y"], Binary(Sub, Var("y"), Var("x"))), UserFn(["x"], Var("x")))
  == Error(WrongArgumentCount(want: 2, got: 1))
\end{lstlisting}

\begin{solution}
\begin{lstlisting}
pub fn compose(f: UserFn, g: UserFn) -> Result(UserFn, EvalError) {
  case f.params {
    [x] -> Ok(UserFn(g.params, substitute(f.body, x, g.body)))
    params -> Error(WrongArgumentCount(want: list.length(params), got: 1))
  }
}
\end{lstlisting}
\end{solution}

\begin{questionmeta}
%   The composed function takes whatever arguments \texttt{g} takes, so its parameters are \lstinline{g.params}. Its body should compute \texttt{f}'s body where \texttt{f}'s parameter is the value of \texttt{g}'s body, so we substitute the whole expression \lstinline{g.body} for \texttt{f}'s parameter. The pattern \lstinline{[x]} matches only a list with exactly one element, which both checks that \texttt{f} has one parameter and gives us its name. In the error, \texttt{f} wants \lstinline{list.length(params)} arguments but would only ever get the 1 value that \texttt{g} produces.

%   Using the same parameter name in both functions is fine. Composing \lstinline|fn(x) { (x * 2) }| with \lstinline|fn(x) { (x + 3) }| gives \lstinline|fn(x) { ((x + 3) * 2) }|, because \texttt{substitute} only makes one pass over \texttt{f}'s body and never looks inside the replacement it puts in.

%   Also note that \texttt{g}'s body is copied into every place \texttt{f} uses its parameter, so in the \texttt{(x * x)} example \lstinline{(y - x)} is evaluated twice. This gives the right answer because evaluating an expression has no side effects, but it does duplicate work.
% \end{questionmeta}
\end{blocksection}